% Transcription — fonds Alexandre Grothendieck, shelfmark 6, batch 3 (pages 41-60).
% Established from the facsimile archives/batches/6/batch-03.pdf, cut from a
% copy of 6.pdf fetched through the project's own relay
% (https://grothendieck-relay.onrender.com/source/6.pdf) rather than from
% grothendieck.umontpellier.fr directly; per the skill transcribe-grothendieck.
% The batch file carries no cover sheet: PDF page k is the archivists' page
% 40+k. Their pencil numbers bottom-left run eleven ahead of the facsimile
% throughout this batch (p. 41 pencilled 52, p. 42 → 53, ..., p. 60 → 71),
% continuing the shift that begins at the end of batch 2 (p. 39 pencilled
% 50); \page{} follows the facsimile, as batch 2 does, and the pencil number
% is not repeated page by page.
%
% Pass: Opus 5.5 (claude-opus-5-5), 2026-09-26 — first pass, unchecked against the pages by a human.
%
% The folder sits in the inventory group "4-9" « Cristaux »; it has its own
% date in the catalogue, which \dating{} carries verbatim, with no group
% sentence.
%
% What the batch contains, and the authorship decisions. The letter to Tate
% ended in batch 2 (p. 31). Everything transcribed here is HIS OWN
% manuscript, in ink, in his hand (his shorthand « tt », « t.m. »,
% « loc. », « kif-kif », « NB », his struck-and-rewritten drafting, the same
% hand as the covers): transcribed in full as his mathematics.
% - Pages 42 and 44: continuation of the run under the cover « Questions
%   diverses en théorie de Tate-Hodge » (batch 2, p. 39; p. 40 there is the
%   first sheet), written, like p. 40, on the blank backs of a foreign
%   typescript (show-through visible). Hodge-Tate structures on V = M ⊗ C
%   (V^{10}, V^{01}, action of π, lattice M); a conjectural Hodge-Tate
%   decomposition of H^i_B(X_C, C) and the comparison ρ^{pq} with
%   H^q(X, Ω^p), with cycle classes P_B(2i) ↦ P_DR(2i).
% - Pages 41, 43, 45: the other sides of those leaves — pages 15, 23, 27 of
%   SOMEONE ELSE'S mimeographed typescript on Banach algebras and the
%   holomorphic functional calculus (numbered « n° 357 », « n° 380 »),
%   unrelated to the folder and carrying no mark of his. Skipped without
%   description (brief and skill: an unrelated verso).
% - Page 46: cover in his hand, « Question d'isogénéité de Barsotti ;
%   impossibilité de plonger un groupe formel p-divisible dans une v.A. sur
%   un k donné (non alg. clos) ». No \page; its words are the heading.
% - Pages 47-50: « Questions sur V.A. et groupes formels »: 1) Barsotti's
%   isogeny question for supersingular abelian varieties with isogenous
%   formal groups, Manin, the counterexample for non-supersingular surfaces
%   with formal group Ĝ_m + G_{1,1}; 2) twisting the formal group of a
%   supersingular elliptic curve over a finite field by a unit θ of the
%   quaternion order, and Weil's obstruction on det φ'. P. 50: three lines,
%   the rest blank.
% - Page 51: cover in his hand, « [Prolongements infinitésimaux d'extension
%   de groupes de B.T. :] Conséquences pour structure infinitésimale des
%   schémas modulaires ». No \page.
% - Pages 52-56: under that cover; p. 52 opens mid-sentence (the previous
%   sheet is not in the folder at this place). Formal moduli of extensions
%   of M by N (N of multiplicative type): M_{ξ_0} a formal torus with
%   reduction T_p(M_0)^∨ ⊗ N_0, M_ξ a torsor under it; compatible pairs of
%   extensions as a kernel of α−β; application to ordinary polarised abelian
%   varieties: local complete intersection, smoothness of the modular scheme
%   in char. p for degree prime to p, dimension g(g+1)/2.
% - Page 57: blank. Skipped.
% - Page 58-59: « Compléments sur les groupes de Witt » (his title,
%   underlined): Lemma 1 on Ext(G_a, G_a), Lemma 2 on Ext(H, W_n), Cor. 3
%   Ext^1(H, W) = 0, Prop. 4 for G annihilated by V^n, and definition 5
%   D(G) = Hom(G, W) with F, V (formula (5.1)). P. 59 is struck through by
%   two long diagonal strokes; transcribed, the cancellation noted once.
% - Page 60: « Groupes finis sur base de car. p>0 » / « Groupes finis
%   étales » (his titles): the crystal D_*(G) for G finite étale, the
%   contravariant D^*(G), F and V, the case G flat over Λ_n (Barsotti-Tate
%   group truncated), anti-equivalence with F-V-crystals. Continues into
%   batch 4.
%
% Pages skipped (no \page): 41, 43, 45 (foreign typescript, unrelated), 46
% and 51 (covers, carried as headings), 57 (blank). Pages summarised: none;
% no page given only as a note.
%
% Notation fixed once for the batch: his round script M (moduli) as
% \mathcal{M}, script N as \mathcal{N}; bold G for his G with a hook
% (formal group) left as plain G; Z_p, Q_ℓ, F_p, C in \mathbf; W underlined
% with an arrow (the ind-group lim W_n) as \underline{W}; Hom, Ext
% underlined (sheaf versions) as \underline{\mathrm{Hom}},
% \underline{\mathrm{Ext}}; H_B for his « H_B » (Betti), H_DR, P_B, P_DR;
% Ĝ_m for his « Ĝm », G_{1,1} as written.
%
% Readings to check first: the whole of pp. 52-56 (fast, pale hand: much
% connective prose is \ill); the margin note of p. 56 (turned);
% « hypersingulières » throughout pp. 47-49; « de dim 2 au plus » p. 47
% (one would expect « 1 »); the matrix at the foot of p. 42; « formelle »
% and the struck block of p. 49; the struck strip at the foot of p. 48.

\documentclass[11pt,a4paper]{article}
% Le préambule commun à toutes les transcriptions du fonds Grothendieck.
%
% Il tient en une idée : **l'incertitude se compose, elle ne se résout pas.**
% Une transcription qui lisse un mot douteux a détruit la seule chose pour
% laquelle on la faisait — un lecteur doit pouvoir savoir, à chaque endroit,
% ce qui était sur le feuillet et ce qui a été ajouté. Les six macros qui
% suivent sont donc l'appareil critique, et non de la décoration.
%
% Les mêmes commandes sont reconnues par `scripts/render.mjs`, qui produit la
% vue de lecture du site. Une macro ajoutée ici doit l'être aussi là-bas,
% faute de quoi le rendu échouera bruyamment — ce qui est voulu : un rendu
% silencieusement faux est pire qu'un rendu absent.

\ProvidesPackage{grothendieck}[2026/08/08 Transcriptions du fonds Grothendieck]

\usepackage{amsmath,amssymb,amsthm}
\usepackage{mathrsfs}
\usepackage{stmaryrd}
% Les diagrammes commutatifs sont partout dans ce fonds : c'est la langue
% même de la géométrie algébrique de Grothendieck, et une transcription qui
% les rendrait en prose ne transcrirait rien.
\usepackage{tikz-cd}
% Français par défaut — c'est la langue des feuillets — mais l'anglais est
% chargé aussi : la traduction et le résumé sont des documents anglais, et la
% typographie française (espaces avant les deux-points) y serait une faute.
% Ces documents basculent par \selectlanguage{english} en tête de corps.
\usepackage[english,french]{babel}
\usepackage{fontspec}
\usepackage{geometry}
\geometry{a4paper,margin=25mm,marginparwidth=28mm}
\usepackage{marginnote}
\usepackage{xcolor}
% ulem, not soul. `soul`'s \st and \ul are built on a character-by-character
% scan that XeLaTeX's font machinery breaks: the document fails with an
% undefined control sequence rather than degrading. ulem does the same two jobs
% and is Unicode-engine safe. `normalem` stops it hijacking \emph, which the
% transcriptions use for Grothendieck's own emphasis.
\usepackage[normalem]{ulem}
\usepackage{hyperref}
\hypersetup{colorlinks=true,linkcolor=black,urlcolor=black,pdfborder={0 0 0}}

% --- Métadonnées du lot -----------------------------------------------------
% Reprises telles quelles par le script de rendu, qui les affiche en tête de
% la vue de lecture. Elles fixent ce que le fichier prétend couvrir : sans
% elles, une transcription flotte sans point d'attache dans un fonds de seize
% mille feuillets.

\newcommand{\folder}[1]{\gdef\grothFolder{#1}}
\newcommand{\batch}[1]{\gdef\grothBatch{#1}}
\newcommand{\leaves}[2]{\gdef\grothFirst{#1}\gdef\grothLast{#2}}
\newcommand{\foldertitle}[1]{\gdef\grothFoldertitle{#1}}
\gdef\grothFolder{?}\gdef\grothBatch{?}\gdef\grothFirst{?}\gdef\grothLast{?}\gdef\grothFoldertitle{}

% --- Le résumé --------------------------------------------------------------
%
% Il ouvre la lecture modernisée, sous le titre « Résumé », et il est composé
% en petit. Ce n'est pas un ornement : c'est le seul passage du document qui
% ne soit pas des mathématiques, et un lecteur doit pouvoir le reconnaître
% avant de l'avoir lu — puis le sauter s'il connaît déjà le sujet, ou s'y
% arrêter s'il ne le connaît pas. Le filet qui le suit marque la reprise du
% corps.

\newenvironment{resume}
  {\small\setlength{\parskip}{0.45em}}
  {\par\medskip\hrule\medskip}

% --- Mots-clés ---------------------------------------------------------------
%
% En anglais, en clôture du résumé de la lecture modernisée : le vocabulaire
% moderne sous lequel chercher ce que les feuillets construisent. Le manifeste
% du site (`npm run manifest`) les extrait pour étiqueter la cote entière —
% cette ligne est la source unique des tags, il n'y a pas de fichier à côté.
\newcommand{\keywords}[1]{\par\smallskip\noindent
  {\footnotesize\textsc{Keywords} --- \textit{#1}}\par}

% --- L'appareil critique ----------------------------------------------------

% Le numéro de feuillet, en marge. C'est l'ancre qui permet de régler un
% désaccord sur une formule en regardant *un* feuillet plutôt qu'une chemise
% de six cents. Le site s'en sert aussi pour tourner les pages du fac-similé
% quand on fait défiler la transcription.
\newcommand{\leaf}[1]{%
  \marginnote{\footnotesize\color{gray!70}\textsf{#1}}%
  \label{leaf:#1}\ignorespaces}

% Illisible : on ne devine pas. Un mot inventé qui se lit comme les autres est
% le pire résultat possible.
\newcommand{\ill}{\textcolor{red!60!black}{[\,\ldots\,]}}

% Lecture douteuse : le mot est donné, et signalé comme incertain.
\newcommand{\uncertain}[1]{\uline{#1}}

% Ajout de l'éditeur — une lettre manquante, un mot sous-entendu.
\newcommand{\add}[1]{\textcolor{blue!50!black}{[#1]}}

% Biffé par Grothendieck lui-même. Ce qu'il a rayé fait partie du document :
% une rature dit souvent où la pensée a changé de direction.
\newcommand{\struck}[1]{\sout{#1}}

% Note du transcripteur, sur l'état matériel du feuillet.
\newcommand{\note}[1]{\par\smallskip{\small\color{gray!60!black}\textit{[#1]}}\par\smallskip}

% Note marginale de Grothendieck, distincte du corps du texte.
\newcommand{\marginal}[1]{\par\smallskip\noindent\hspace{1em}%
  {\small\color{gray!60!black}#1}\par\smallskip}

% --- Titre ------------------------------------------------------------------

\AtBeginDocument{%
  \begin{center}
    {\large\bfseries Fonds Alexandre Grothendieck --- cote n\textsuperscript{o}~\grothFolder}\\[2pt]
    {\itshape\grothFoldertitle}\\[4pt]
    {\small feuillets \grothFirst--\grothLast\ (lot \grothBatch)}\\[2pt]
    {\footnotesize\color{gray!60!black}Transcription établie d'après le fac-similé numérisé
     par l'Université de Montpellier.\\
     Les lectures incertaines sont soulignées, les illisibles notées [\,\ldots\,],
     les ajouts entre crochets.}
  \end{center}
  \vspace{1em}}

\endinput


\folder{6}
\batch{3}
\pages{41}{60}
\dating{1966-[à partir de 1970]}
\watermark{Édition de démonstration}
\foldertitle{Lettre Tate (mai 66) (Cristaux) : lettre (1966), tapuscrit (s.d.), notes manuscrites (s.d.)}

\begin{document}

\page{42}
\note{suite des « Questions diverses en théorie de Tate-Hodge » (lot
précédent, p.~40), à l'encre, au dos d'une page dactylographiée étrangère
au dossier. Les trois premiers points sont réunis par une accolade dans la
marge de gauche.}

1) $V=V^{10}+V^{01}$ décomposition de vectoriels sur $\mathbf{C}$.

\struck{1') $V^{10}\simeq V_0^{10}\otimes\mathbf{C}$,
$V^{01}\simeq\check{V}^{1}\otimes_K\mathbf{C}[-1]$}\note{ligne biffée d'un
trait ondulé ; lecture douteuse des indices.}

2) Opération de $\pi$ sur $V$, semi-linéaire \uncertain{relativement} à
\uncertain{l'opération} de $\pi$ sur $\mathbf{C}$.

3) $V\simeq M\otimes_{\mathbf{Z}_p}\mathbf{C}$, $M$ un \struck{vectoriel}
\add{\uncertain{module libre de t.f.}} sur $\mathbf{Z}_p$.

\emph{Relations} \emph{1) et 2)}
\[
  V^{10}=V^{\pi}\otimes_K\mathbf{C},\qquad
  V^{01}=\bigl(V\otimes\mathbf{C}[1]\bigr)^{\pi}\otimes_K\mathbf{C}[-1]
\]
\marginal{NB Structures 1, 2 \ill{} \ill{} \ill{} $V^{\pi}$,
$V'=(V\otimes\mathbf{C}(1))^{\pi}$. NB.\ $2\Rightarrow 1$}\note{encadré à
droite des formules, séparé par un trait vertical ; « $2\Rightarrow 1$ »
est écrit avec une double flèche épaisse.}

\emph{2) et 3)} \quad $M$ stable par $\pi$ \quad
[NB.\ structures 2, 3 \ill{} \ill{} \ill{} $M$ avec \ill{} \ill{} de $\pi$
dessus \struck{NB.\ \ill{}} [i.e.\ le groupe $p$-divisible sur $K$]]

\emph{3 et 1} \quad ??? \quad \struck{Les \ill{} \ill{} $M$ \ill{} \ill{}}

\emph{NB} Le corps $K$ intervient dans 2), \ill{} \ill{}
\struck{implicite}, mais pas dans 1) ni 3), qui \uncertain{traduisent} des
structures intrinsèques à la donnée du groupe $p$-divisible sur
$\widehat{S}=\widehat{\Omega}=\mathbf{C}$.

Regardons le \add{plus petit} groupe \struck{alg.}\ algébrique
d'automorphismes \add{$G$} de $M\otimes\mathbf{Q}_\ell$
\struck{\ill{}} tel que $G_{\mathbf{C}}$ contienne les opérations
$\begin{pmatrix}\lambda&0\\0&\xi\end{pmatrix}$ sur la décomposition de $V$
en $V^{10}+V^{01}$.\note{la matrice est de lecture douteuse (on lit
$\lambda$, un $0$, puis $\xi$, avec des points au-dessus) ; « $G$ » est
écrit au-dessus de la barre de $M\otimes\mathbf{Q}_\ell$.}

\page{44}
\emph{Conjecture} 1) $H^i_B(X_{\mathbf{C}},\mathbf{C})$ est somme
\add{directe} des sous-espaces engendrés par les \ill{}
\uncertain{$\chi^{p}$-propres}, de façon intrinsèque, i.e.\ \ill{}, pour
$p+q=i$
\[
  H^{p,q}_B(X)_K=H^i_B(X_{\mathbf{C}},\mathbf{C})[p]^{\pi}
\]
on a une décomposition en somme
\[
  \boxed{\;H^i_B(X_{\mathbf{C}},\mathbf{C})\simeq\coprod_{p+q=i}
  H^{p,q}_B(X)_K\otimes_K\mathbf{C}[-p]\;}
\]

2) On a un isom canonique
\[
  \boxed{\;\rho^{pq}\colon H^{pq}_B(X)_K\simeq H^{p,q}(X)
  \overset{\mathrm{d\acute{e}f}}{=}H^q(X,\Omega^p_X)\;}
\]
\note{les deux formules sont encadrées sur la page. Le $\rho^{pq}$ est
récrit sur un autre signe.}

\[
  P_B(2i)\in H^{2i}(X_{\mathbf{C}},\mathbf{Z}_p[i])^{\pi}
  \subset H^{2i}_B(X_{\mathbf{C}},\mathbf{C})[i]^{\pi}=H^{ii}_B(X)_K,
  \qquad H^{2i}(X_{\mathbf{C}},\mathbf{Z}_p)
\]
\note{la formule est précédée d'un trait horizontal sur toute la largeur.
Le « $\subset H^{2i}(X_{\mathbf{C}},\mathbf{Z}_p)$ » de la ligne du
dessous est placé sous « $H^{2i}_B(X_{\mathbf{C}},\mathbf{C})[i]^{\pi}$ » ;
l'indice du premier $H^{2i}$ est récrit, illisible.}

\struck{$P_B(2i)$} $P_{\mathrm{DR}}(2i)\in H^{i,i}(X)$ ; on veut
\[
  \boxed{\;\rho^{ii}\bigl(P_B(2i)\bigr)=P_{\mathrm{DR}}(2i)\;}
\]

\emph{NB.} On aura \ldots\quad
$P_B(2i)\in H^{ii}_B(X)_{\mathbf{Q}}\cap H^{2i}(X_{\mathbf{C}},\mathbf{Z}_p[i])$
\add{identifié à $H^{i,i}(X)$}\note{l'indice « $\mathbf{Q}$ » de
$H^{ii}_B(X)$ est récrit et se lit mal ; « identifié à $H^{i,i}(X)$ » est
écrit sous la barre qui souligne $H^{ii}_B(X)_{\mathbf{Q}}$.}

\section*{Question d'isogénéité de Barsotti ; impossibilité de plonger un
groupe formel $p$-divisible dans une v.A.\ sur un $k$ donné (non alg.\
clos)}
\note{titre écrit de sa main, en haut à droite d'une feuille par ailleurs
blanche (p.~46 des archivistes), qui sert de chemise aux feuilles qui
suivent ; un mot biffé avant « impossibilité ».}

\page{47}
\emph{Questions sur} \emph{V.A.\ et groupes formels}.

1) (Barsotti). Soient $A,B$ deux VA sur $k$ alg.\ clos (car $p>0$)
« \uncertain{hypersingulières} » i.e.\ sans pts d'ordre $p$. Supposons que
les groupes formels $\overline{A},\overline{B}$ soient isogènes. Alors $A$ et
$B$ sont-elles isogènes ? \add{Si $\dim A=1$, c'est contenu dans
Deuring.}\note{ajout dans l'interligne.} Manin ignore déjà la réponse si
$A,B$ de dim 2 (donc $\overline{A},\overline{B}$ sont isogènes :
$G_{1,1}+G_{1,1}$). La question revient donc à ceci : si $A$ de dim 2 est
\uncertain{hypersingulière}, est-elle isogène à un produit de deux courbes
elliptiques (ou : au carré d'une courbe elliptique) nécessairement
\uncertain{hypersingulière}. Il affirme que c'est vrai en car $p=2$.

Notons que si on abandonne l'hypothèse \uncertain{hypersingulière}, et
qu'on suppose seulement $A,B$ \struck{\ill{}} de dim 2, à groupes formels
isomorphes non toroïdaux, il n'est pas vrai en général que $A$ et $B$
\add{\uncertain{soient}} isogènes \add{$=$ si $A$ et $B$ simples} (en
effet, d'après Manin, les \add{\uncertain{polarisées}} $A$ de dim 2
\struck{(polarisations principales)} de groupe formel
$\widehat{G}_m+G_{1,1}$ forment une sous-famille \add{$N$} \ill{} de
\emph{dim 2} de la famille modulaire \add{$M_2$}. Parmi celles-ci, celles
qui sont \struck{isogènes} \add{non} simples i.e.\ \struck{décomposées}
isogènes à un produit $E\times F$ ($E$ et $F$ courbes elliptiques dont l'une
\add{(soit $F$)} est \uncertain{hypersingulière}, l'autre non), forment une
famille de dim \uncertain{2} au plus $\lbrace$\emph{NB.} il n'y a pas de famille
continue\note{les ajouts sont écrits au-dessus de la ligne ; « dim 2 » est
souligné. On attendrait « dim 1 au plus » : lecture douteuse du chiffre.}

\page{48}
de variétés isogènes à $E\times F$ !$\rbrace$ Plus précisément, les
\uncertain{courbes} isogènes à $E\times F$ \struck{\ill{}} sont -- sur un
corps alg.\ clos -- \add{isomorphes} \struck{isogènes} à un produit
$E'\times F'$. Donc les variétés \uncertain{isogènes} à un produit forment
une sous-famille \ill{} de $E$, image de l'ouvert $M'_1$ de la variété
modulaire $M_1$ des courbes elliptiques, formé de celles qui sont
ordinaires, par les applications $E\mapsto E\times F$, où $F$ parcourt
\ill{} fini de types de courbes ell.\ \uncertain{hypersingulières}. D'où un
ouvert non vide $U$ dans $N$ \struck{\ill{}} formé de variétés simples.
\struck{Notons que si $A$ est une variété abélienne \ill{} \ill{} simple
(i.e.\ qui \ill{} \ill{} \ill{} \ill{} \ill{} \ill{} de dim 1) elle n'est
pas \ill{} \ill{} un produit $E\times F$}\note{passage biffé de traits
horizontaux et barré de plusieurs grandes croix ; lecture très
incomplète.} D'autre part si $A$ \ill{} \uncertain{simple} \ill{} \ill{}
\ill{} la structure formelle envisagée, dans ce \ill{} \ill{} dessus que les
$A'$ isogènes à $A$ forment une famille \emph{discrète}, (\ill{} \ill{}),
de façon précise si $A$ est définissable sur $k_0$, alors $A'$ est
définissable sur une extension \emph{finie} de $k_0$. Si on prend deux
points \ill{} de \struck{\ill{}} $U\subset N$ \struck{l'un \ill{} \ill{},
l'autre} distincts ayant des degrés de transcendance distincts sur le corps
premier $\mathbf{F}_p$ (les degrés possibles sont 0, 1, 2) on trouve des
variétés simples \emph{non isogènes}, de même groupe formel
$\widehat{G}_m\times G_{1,1}$.

\page{49}
2) Soit $k$ un corps, par ex.\ alg.\ clos. Alors en général, un groupe
formel \add{$G$} sur $k$, même \add{une forme de} $G_{1,1}$
\uncertain{(peut-être)}, ne peut se réaliser comme s-gpe formel du groupe
formel associé à une VA définie sur $k$ -- et ceci même si on se permet une
extension finie \struck{de $k$} (ou de t.f.)\ de $k$. Prenons p.ex.\ $k$
corps fini, $E$ courbe elliptique \uncertain{hypersingulière} sur $k$
\add{« \uncertain{à multiplic.\ complexe définie dans $k$} »},
\struck{\ill{}} ($G_0=\widehat{E}$), donc $\pi=\pi_1(k)$ opère
trivialement sur $\mathrm{End}(E)$ et $\mathrm{End}(G_0)$
\add{$\simeq\mathrm{End}(E)\otimes_{\mathbf{Z}}\mathbf{Z}_p=\mathcal{A}$},
qui sont \add{\uncertain{chacun un}} ordre maximal dans des
\struck{\ill{}} algèbres de quaternions sur $\mathbf{Z}$ resp.\
$\mathbf{Z}_p$. On peut tordre \ill{} \ill{} \uncertain{unité} $\theta$ de
$\mathcal{A}$, dans $G_0$ par \ill{} \add{\uncertain{formelle}} \ill{} :
$\varphi'=\theta\varphi$, $\varphi$ l'ancien frob., \struck{le nouveau
frobenius devant \ill{} jouant le rôle de frobenius}. [Mais si on veut que
$G_0$ se plonge dans un $\widehat{A}$, $A$ une VA, il faut par Weil que
les \ill{} de l'équation caractéristique \ill{} de $\varphi'$
\add{$=\theta\varphi$} \struck{\ill{} \ill{} \ill{} \ill{} \ill{} $Z=$ (ce
qui \ill{} \ill{} \ill{} \ill{} \ill{} \ill{} de val.\ abs.\ $q^{1/2}$,
p.ex.\ il faut \ill{} \ill{} \ill{} \ill{}), \ill{} \ill{}} \struck{\ill{}
il faut \ill{} \ill{} (\ill{} \ill{} $q^{1/2}$, \ill{} il faut $\det\varphi'$
$\det\varphi'=\pm\det\varphi$},\note{bloc de quatre lignes biffé de traits
horizontaux et barré de deux obliques ; lecture très incomplète.} ce qui
exige que $\det\varphi'=\det\varphi\,\det\theta$ \struck{\ill{}} soit une
certaine \ill{} de \uncertain{val.\ abs.\ égale}, ce qui implique que
$\det\theta$ doit être \ill{} \ill{} dont \ill{} \ill{} conjugués sont de
\uncertain{val.\ abs.} 1.

\page{50}
Mais bien entendu il est facile de trouver $\Theta$ unité de $\mathcal{A}$
\struck{\ill{} \ill{}} tel que $\det\Theta$ ne soit pas algébrique
\ldots\ldots\note{le $\Theta$ final de la première ligne est d'une autre
encre ; le reste de la page est blanc.}

\section*{[Prolongements infinitésimaux d'extension de groupes de B.T.\ :]
Conséquences pour structure infinitésimale des schémas modulaires}
\note{titre écrit de sa main, en haut à droite d'une feuille par ailleurs
blanche (p.~51 des archivistes) ; les crochets sont les siens. Un mot
biffé avant « modulaires ».}

\page{52}
\note{la page commence au milieu d'une phrase ; la feuille qui précède
n'est pas à cette place dans le dossier. Une grande accolade au crayon
marque la fin du premier alinéa.}
$S^{(n)}\to G=\widehat{S}$, \add{\uncertain{nul sur} $S^{(0)}$,} les hom
se recollent et définissent un hom de prolongements formels
\[
  S^{(\infty)}\longrightarrow G
\]
dont la connaissance \struck{\ill{}} \add{\uncertain{équivaut à}} la
connaissance de l'ext.\ donnée $E$ de $M$ par $N$, non pas sur $S$
lui-même, mais dans le voisinage inf.\ de $S$ sur $R$, et \ill{} \ill{}
\ill{} \ill{} de $S$ sur $R$, et kif-kif après changement de base.

La formation de $S^{(\infty)}\to G$ est compatible avec changement de base
\struck{\ill{}} $S'\to S$. Il \ill{} \ill{} donc de la considérer, dans le
\ill{} des positifs, comme \uncertain{découlant} d'une situation
« \emph{modulaire} ».\note{« modulaire » est souligné deux fois au crayon.}

Supposons \struck{\ill{}} \add{\uncertain{d'abord}} que l'on \ill{}
\add{\ill{}} \ill{}, $R=\operatorname{Spec}\Lambda$, $\Lambda$ anneau local
complet de car.\ résiduelle $p$ (p.ex.\ un \ill{} de Cohen), et qu'on se
donne \ill{} de corps résiduel $k$, $M_0$ et $N_0$ sur $k$. On veut la
variété \add{formelle} \add{$(\mathcal{M}_\xi)$} sur $R$ (considérée comme
schéma formel) modulaire \add{\ill{}} des extensions \ill{} $\xi$ de
$M_0$ par $N_0$. Supposons $N_0$ de t.m.\ \struck{\ill{}} Conditions sur la
variété modulaire formelle pour $\xi_0$, l'extension triviale.
\struck{\ill{}} Alors :\note{les ajouts, au-dessus et au-dessous de la
ligne, sont reliés par un long trait ; lecture douteuse de leur ordre.}

\page{53}
1) $\mathcal{M}_{\xi_0}$ est \supplied{la structure d'}un groupe formel lisse
sur \struck{$R$}, $\operatorname{Specf}\Lambda$, grâce à la structure de
groupe du foncteur qu'il représente. Le groupe formel en question est un
tore \emph{formel}, \struck{il est \ill{} \ill{}} \add{\uncertain{donc
déterminé}} par sa réduction sur $k$, et cette dernière est le
\add{\ill{}} \struck{tore}
$T_p(M_0)^{\vee}\otimes N_0$.

2) $\mathcal{M}_\xi$ est un \struck{fibré principal homogène}
\add{torseur} sous $\mathcal{M}_{\xi_0}$. [qui \ill{} \ill{} \ill{} \ill{}
\ill{} un pseudo-torseur \add{provenant du th.\ de Tate affirmant que
\ill{} \ill{} \ill{} \ill{}} \ill{} des groupes \struck{formels}
$p$-divisibles]

Ne supposons plus \ill{} $N_0$ de t.m., mais \ill{} \uncertain{radiciel}.
On veut préciser la structure de la variété formelle des modulaires pour
$E_0$ [NB.\ $E_0$ détermine uniquement $M_0$, $N_0$ et la structure
d'extension !]\ On a un hom naturel
\[
  X=\mathcal{M}_{E_0}\longrightarrow\mathcal{M}_{N_0}=Y .
\]
Ce dernier fait de $\mathcal{M}_{E_0}$ un \struck{fibré} torseur
\struck{principal homogène} sous
$\mathcal{N}_Y\otimes_{\mathbf{Z}_p}T_p(E_0)^{\vee}_Y$. Donc la structure
de \struck{$X$} $\mathcal{M}_{E_0}$ est connue si celle de
$\mathcal{M}_{N_0}$. [Dans le cas $N_0$ de t.m., $Y$ est réduit à un
point].\note{dans « $\mathcal{M}_{E_0}$ » de la formule, l'indice est
récrit sur un autre ; de même pour $\mathcal{M}_{N_0}$ à la fin.}

\page{54}
Revenons au \emph{pb} des extensions de $M$ par $N$, \emph{triviales} sur
$S_0$. Soient \ill{} \ill{} des $M'$, $N'$ et \ill{} \ill{} le \ill{} pb.
Donnons-nous de plus des hom $M\xrightarrow{u}M'$, $N\xrightarrow{v}N'$, et
donnons-nous des systèmes d'extensions $E$ de $M$ par $N$, $E'$ de $M'$ par
$N'$, \add{\uncertain{triviales} sur $S_0$} \struck{\ill{}} tels que
$\exists$ \struck{\ill{}} hom $E\to E'$ compatible avec $u,v$ [on
\ill{} pas \ill{} unique].

On a \ill{} $\mathrm{Ext}^1(M,N)\xrightarrow{v_*}\mathrm{Ext}^1(M,N')$ et
$\mathrm{Ext}^1(M',N')\xrightarrow{u^*}\mathrm{Ext}^1(M,N')$, des hom de
groupes, \struck{\ill{} \ill{} \ill{}} et les couples $(\xi,\xi')$
cherchés sont ceux pour lesquels $v_*(\xi)=u^*(\xi')$. Or l'hom $v_*$ est
« représenté » par l'hom de \emph{tores}
\[
  T_p(M)^{\vee}\otimes N\xrightarrow{\ \alpha\ }T_p(M)^{\vee}\otimes N'
\]
défini par $v$, et de même $u^*$ par
\[
  T_p(M')^{\vee}\otimes N'\xrightarrow{\ \beta\ }T_p(M)^{\vee}\otimes N'
\]
défini par $u$. Ceci dit, si \uncertain{on considère} le \emph{noyau}
\struck{de \ill{}} \add{i.e.\ de $\alpha-\beta$,} $G\times
G'\rightrightarrows H$, \add{\ill{} \ill{}} ce dernier représente
infinitésimalement le système des\note{« $v_*$ », « $u^*$ » sont écrits
au-dessus des flèches ; le « $\alpha$ » et le « $\beta$ » aussi. Le premier
$T_p$ est récrit sur un autre signe.}

\page{55}
couples $(\xi,\xi')$, avec les propriétés de compatibilité voulues.

Considérons en \ill{} \ill{} \ill{} situation de modules formels
\add{(\uncertain{non nécessairement} \ill{})} avec $N,N'$ de t.m., on
\ill{} \ill{} que \ill{} \ill{} \ill{} d'extensions triviales) la variété
formelle des modules pour le \ill{} \ill{} est un groupe formel \add{de
t.m.}\ sur $\Lambda$ \struck{(NB.\ \ill{} \ill{} \ill{} !)}, noyau du hom
$\alpha-\beta\colon\mathcal{M}\to\mathcal{N}$ \struck{\ill{}} [qui \ill{}
un hom \emph{de groupes}]. Cela s'applique notamment au \ill{} formel des
modules pour une V.A.\ \emph{polarisée}. Quand $\Lambda$ est le spectre
\ill{} d'un corps de car.\ $p$, cela implique \ill{} \ill{} que l'anneau
local obtenu est \struck{\ill{}} une intersection C.M.\ (et
\struck{\ill{}} \ill{} \ill{} complète, \ill{} Gorenstein). D'où le fait
que les schémas \struck{\ill{}} modulaires pour les V.A.\ polarisées
\add{\emph{ordinaires}} (\ill{} \ill{}, \ill{} \ill{}\note{« de t.m. »,
« ordinaires » et la parenthèse de la deuxième ligne sont écrits
au-dessus de la ligne et appelés par un trait. Le nom de l'hom
« $\alpha-\beta$ » et ses deux membres sont de lecture douteuse.}

\page{56}
polarisation donnée, et rigidification de Jacobi \add{\uncertain{d'échelon}}
(\uncertain{donné}), en car.\ $p$, est \ill{} une intersection complète, et
de plus est \emph{loc.} \emph{irréductible}, plus précisément
$\mathcal{M}_{\mathrm{red}}$ \struck{\ill{}} \add{\uncertain{lisse sur
$k$}}. De plus, \struck{\ill{}} \struck{dimension de la \ill{} correspondant
à une \ill{} \ill{} \ill{} \ill{} \ill{} et utilisant la}
\marginal{\uncertain{dimensionnalité} \uncertain{ridicule}, on le voit
\uncertain{bien} directement en car.~$p$ !}\note{note écrite en travers
dans la marge de gauche, à hauteur du passage biffé ; lecture douteuse.}
Utilisant \ill{} \ill{} que \emph{la dimension des schémas modulaires}
\add{\uncertain{ordinaires}} \struck{est} \emph{bien}
$\frac{g(g+1)}{2}$ \emph{en chaque point}. Donc si la polarisation est de
degré premier à $p$ et \ill{} \ill{} en car.\ zéro, le schéma modulaire en
car.\ $p$ est lisse. (NB.\ il ne s'agit que de VA
\emph{ordinaires})\note{la ligne « dimension de la \ldots\ correspondant à
une \ldots » est barrée d'un trait qui court sur toute la largeur ; le
reste de la page est blanc.}

\section*{Compléments sur les groupes de Witt}
\note{titre souligné, de sa main, en tête de la p.~58 ; la p.~57 est
blanche.}

\page{58}
\emph{Lemme} 1 \quad Soit $S=\operatorname{Spec}(A)$ de car.\ $p$.
\struck{Posons $A[\mathbf{F}]$, $\mathbf{F}_p[$}

$\underline{\mathrm{Ext}}^1(G_a,G_a)$ est le \struck{$A[\mathbf{F}]$-module}
$A$-module \add{\uncertain{à gauche}} \uncertain{libre} engendré par les
$\mathbf{F}^i e$, \add{i.e.\ le $A[\mathbf{F}]$-module libre engendré par
$e$,} (où $e$ est la classe de $W_2$ considéré comme extension de $G_a$ par
$G_a$.\note{le « $A[\mathbf{F}]$-module \ldots\ par $e$ » est écrit
au-dessus de la ligne ; « à gauche », entouré, est appelé d'en bas par un
trait. L'énoncé est marqué d'un trait vertical dans la marge.}

\emph{Dém} On peut supposer, par \uncertain{\ill{}}, que
$A=\mathbf{F}_p$. Mais c'est fait dans Demazure-Gabriel, II~§3 4.6

\emph{Lemme} 2 \quad Soit $H$ un sous-schéma \struck{\ill{} \ill{} \ill{}
\ill{} \ill{} de $G_a$ tel \ill{} \ill{} $G/H$ \ill{}} \add{\ill{}} de
$G_a$, tel \struck{\ill{}} que \ill{} \ill{} $H\simeq G_a$, où $G_a/H$ est
loc.\ isom.\ à $G_a$ (groupe \uncertain{additif}). Alors \ldots
\begin{itemize}
  \item[a)] $\underline{\mathrm{Ext}}^1(H,W_n)\xrightarrow{R^{n-1}*}
  \mathrm{Ext}^1(H,W_1)$ est bijective,

  \struck{$\mathrm{Ext}^1(H,W_{n-1})\to\mathrm{Ext}^1(H,W_n)$ est nulle
  \ill{}}

  \item[a$_1$)] $\mathrm{Ext}^1(H,W_{n-1})\to\mathrm{Ext}^1(H,W_n)$ est
  \uncertain{nulle}, i.e.\ $\mathrm{Hom}(H,G_a)\to\mathrm{Ext}^1(H,W_{n-1})$
  est surjectif
  \item[a$_2$)] $\mathrm{Ext}^1(H,W_1)\to\mathrm{Ext}^1(H,W_{n-1})$ est
  nulle
  \item[b)] $\mathrm{Ext}^1(G_a,W_n)\to\mathrm{Ext}^1(H,W_n)$ surjective
  (i.e.\ $\mathrm{Ext}^1(H,W_n)\to\mathrm{Ext}^2(G_a/H,W_n)$ est nulle),
  \item[c)] $\underline{\mathrm{Ext}}^1(G_a,W_n)$ est un
  $A[\mathbf{F}]$-module \add{\uncertain{à droite}, engendré par la}
  \struck{\ill{} \ill{} \ill{}} classe \struck{\ill{}} de $W_{n+1}$.
\end{itemize}
\note{entre a) et a$_1$), une ligne biffée et un petit schéma
« $\mathrm{Ext}^1 \downarrow V$ » de lecture douteuse, qui renvoie à la
flèche de a$_1$). Dans a), l'exposant de $R$ se lit « $n-1$ ».}

\emph{Dém} \emph{n=1} Alors a) est trivial, b) résulte de
$\mathrm{Ext}^2(G_a/H,W_1)=0$ (qui résulte, par dévissage \struck{\ill{}},
de \add{loc.} $\mathrm{Ext}^2(G_a,G_a)=0$) et c) n'est autre que le lemme 1.
\add{\uncertain{V} §1 2.2.}\ \add{\emph{n général} :} On procède comme dans
\struck{\ill{}} loc.\ cit., par récurrence sur $n$.

\page{59}
\note{toute la page est barrée de deux longs traits obliques.}
\emph{Corollaire} 3 \quad Si $H$ est comme précédemment, \ldots\
$\underline{\mathrm{Ext}}^1(H,\underline{W})=0$
($\underline{W}=\varinjlim W_n$).\marginal{\emph{Question} Si $S$ est
\ill{}, est-ce que $\mathrm{Ext}^1(H,\underline{W})=0$, i.e.\
$H^1(S,\underline{\mathrm{Hom}}(H,\underline{W}))=0$ ??}

\emph{Non} déjà pour $H=\alpha_p$ si $S$ \add{\uncertain{parfait,} \ill{}
\ill{} \ill{}} non \uncertain{parfait}.

\emph{Proposition} 4 \quad Soit $G$ un groupe algébrique \struck{\ill{}}
\add{\ill{}} \ill{} \ill{} \ill{} \add{(par \uncertain{ex.} parfait)} tel
que $G$ soit annulé par $V^n$. Alors
\begin{itemize}
  \item[a)] $\underline{\mathrm{Hom}}_{\mathrm{gr}}(G,W_n)$ est un schéma en
  groupes \struck{\ill{}} \add{\uncertain{annulé}} \struck{\ill{} \ill{}}
  \struck{\ill{} $G^0$} (\uncertain{et} \uncertain{isomorphe} :
  $\mathbf{E}^r_k$ \ill{} $r$ \struck{\ill{} \ill{}}, $\mathrm{rg}\,G=p^r$)
  \item[b)] $\underline{\mathrm{Ext}}^1_{\mathrm{gr}}(G,\underline{W})=0$
\end{itemize}

Le b) provient de cor.~3, compte tenu que \ill{} \ill{} un groupe \ill{} :
$\alpha_p$.

Pour prouver a), \ill{} procède par récurrence sur la longueur $r$
d'\add{\uncertain{une}} suite de composition. On a donc
\[
  0\to G'\to G\to\alpha_p\to 0
\]
d'où une suite exacte
\[
  0\to\underline{\mathrm{Hom}}(\alpha_p,\underline{W})\to
  \underline{\mathrm{Hom}}(G,\underline{W})\to
  \underline{\mathrm{Hom}}(G',\underline{W})\to 0
\]
compte tenu que $\underline{\mathrm{Ext}}^1(\alpha_p,\underline{W})=0$. On
peut \ill{} \ill{} \ill{} \ill{} $\underline{W}$ par $W_n$. \ill{} \ill{},
on applique l'hyp.\ de récurrence : $\underline{\mathrm{Hom}}(G',W_n)$
\ill{} \ill{} \ill{} que
$\underline{\mathrm{Hom}}(\alpha_p,W_n)\simeq G_a$. On gagne !

5. Soit $\mathbb{D}(G)=\underline{\mathrm{Hom}}(G,\underline{W})$. C'est un
$W_\infty$-Module, \ill{} \ill{} des \ill{} $V$ et $\mathbf{F}$
satisfaisant les conditions habituelles (\uncertain{provenant} des
relations entre $V$, $\mathbf{F}$ et les $\lambda\in W_\infty(S)$,
\uncertain{opérant} dans $\underline{W}$)
\[
  \mathbf{F}\lambda=\lambda^{(p)}\mathbf{F},\qquad
  \lambda V=V\lambda^{p},\qquad
  \mathbf{F}V=V\mathbf{F}=p.\mathrm{id}
  \tag{5.1}
\]
\note{la question est écrite à droite du corollaire, au-dessus de la ligne
suivante ; dans a), « $\mathbf{E}^r_k$ » (un $E$ à double trait) et la
parenthèse sont de lecture douteuse, une partie de la ligne étant biffée.}

\section*{Groupes finis sur base de car.\ $p>0$}
\note{titre de sa main en tête de la p.~60, non souligné ; le suivant est
souligné.}

\page{60}
\emph{Groupes finis étales}

Soit $G$ fini étale annulé par $p^n$. Pour tout épaississement $S'$ de $S$,
$G$ s'étend de façon unique en un groupe fini étale $G'$ sur $S'$,
également tué par $p^n$. Soit
$\mathcal{M}_{S'}=G'\otimes_{\mathbf{Z}_p/p^n\mathbf{Z}}\underline{O}_{S'}
=G'\otimes_{\mathbf{Z}_p/p^n\mathbf{Z}}(\underline{O}_{S'}/p^n\underline{O}_{S'})$.
Pour $S'$ variable, on obtient un cristal $\mathbb{D}_*(G)$. L'hom
$G\xrightarrow{f}G^{(p)}$ est ici un isom.\ et définit un isom
\[
  F_*\colon\mathbb{D}_*(G)\longrightarrow\mathbb{D}_*(G)^{(p)}
\]
qui est un isomorphisme. Par dualité, posant
\[
  \check{G}=\underline{\mathrm{Hom}}_{\mathrm{gr}}\bigl(G,(\mathbf{Q}_p/\mathbf{Z}_p)_S\bigr),
  \qquad \mathbb{D}^*(G)=\mathbb{D}_*(\check{G}),
\]
on trouve un \supplied{contra-}foncteur $G\mapsto\mathbb{D}^*(G)$, et
\[
  F^*\colon\mathbb{D}^*(G)^{(p)}\xrightarrow{\ \sim\ }\mathbb{D}^*(G).
\]
$V^*$ est uniquement défini par \struck{\ill{}}
$F^*V^*=p\,\mathrm{id}_{\mathbb{D}^*(G)}$.

\struck{Dans la catégorie des cristaux de modules}\note{ligne biffée, d'une
écriture posée.}

\emph{Cas} $G$ \emph{plat sur} $\Lambda_n$ (i.e.\ groupe de B.T.\ tronqué
d'échelon \ldots)

\quad a) Alors $\mathbb{D}^*(G)$ est un cristal de modules loc.\ libres de
t.f.\ rel.\ à $\Lambda_n$, et avec $F$ \struck{tel que} isomorphisme
($\Rightarrow V$ de façon unique) et on obtient une \supplied{anti-}équivalence
de la catégorie des \struck{\ill{}} $G$ et des $F$-cristaux
précédents.\marginal{\uncertain{vérifier}}\note{la note marginale est
écrite de biais à gauche d'une accolade qui embrasse l'alinéa a).}

\emph{Cas général} (\supplied{Anti-}équivalence entre la catégorie des $G$, et
celle des \supplied{$F$-}cristaux (i.e.\ $F$-$V$-cristaux) sur $S$ qui loc.\
(ét) sont des sommes de cristaux du type précédent (avec $n$ variable)).

De ceci résulte que si $G\leftrightarrow\mathcal{M}$, \ldots,
\[
  G(S)\underset{\mathrm{unpd.}}{\simeq}
  \mathrm{Hom}_{\mathrm{gr}}(\Lambda_{n\,S},G)\simeq
  \mathrm{Hom}_{F\text{-}V\text{-}\mathrm{cris}}(\mathcal{M},
  \underline{O}_{S_n\,\mathrm{cris}})\simeq\mathrm{Hom}_{F\text{-}\mathrm{cris}}(\ldots
\]
\note{la formule court jusqu'au bord de la feuille et se poursuit au lot
suivant ; l'indice sous le premier $\simeq$ est de lecture douteuse.}

\end{document}
